\documentclass[10pt,a4paper]{amsart}
\usepackage[margin=19mm]{geometry}
\usepackage{amsmath,amssymb,mathtools}
\usepackage[scr=rsfs,cal=euler]{mathalfa}
\usepackage{xfrac}
\usepackage[unicode,hidelinks,bookmarks=false]{hyperref}
\newcommand{\ie}{i.e.\ }
\newcommand{\cf}{cf.\ }
\newcommand{\resp}{respectively\ }
\newcommand{\Subsection}{\subsection}
\providecommand{\nobreakdash}{\nobreak\hbox{-}}
\setlength{\emergencystretch}{2em}
\multlinegap0pt
\usepackage{amssymb}
\usepackage{mathtools}
\usepackage{bm}
\newtheorem{lemma}{Lemma}
\title{Microlocal analysis of double fibration transforms with conjugate points}
\author{Hiroyuki Chihara}
\date{Source: arXiv:2412.14520v3 (CC BY 4.0)}
\begin{document}
\maketitle

\noindent\emph{Source and licence.} Microlocal analysis of double fibration transforms with conjugate points. Version arXiv:2412.14520v3 (CC BY 4.0), distributed by arXiv under the Creative Commons Attribution 4.0 International licence, http://creativecommons.org/licenses/by/4.0/, as recorded in the arXiv licence metadata for this exact version and shown on the abstract page https://arxiv.org/abs/2412.14520v3. Full licence notice: \texttt{licenses/arxiv-2412.14520-CC-BY-4.0.txt}.\par\smallskip

\noindent\textit{Editorial scope.} Counted unit: the article's lemma theorem:submanifold (source0.tex lines 1268--1275). The article's complete original proof of it is delivered with it (source0.tex lines 1278--1388, one \texttt{proof} environment of 111 lines). No mathematics is written, repaired or re-derived: the statement and the proof are the article's own lines, in the article's own notation, and no retained line is substituted or re-flowed.  No further source-local context is needed: every cross-reference of every retained line resolves inside the two delivered spans.  Nothing was omitted from any retained span, so the package carries no omission record.  No retained line contains a citation, so the wrapper carries no bibliography and no bibliography entry.  No retained line contains a figure environment, an image-inclusion command or drawing code, so no image is delivered and the package declares no asset dependency.  Editorial formatting only, in addition: an AMS \texttt{amsart} preamble that restates the article's own declarations and macros used by the retained lines, namely the environments \texttt{lemma} on separate counters, together with the standard packages those lines need.  The retained environments print in this document as Lemma 1 in order of appearance, whereas the article prints the same items with the numbering of its own full document.  The complete native source of the article as distributed in the arXiv e-print for this exact version is bundled unchanged as \texttt{sources/170379/source0.tex}.
\begin{lemma}
\label{theorem:submanifold} 
Suppose that $Z$ is a double fibration and $N \geqq 2n^{\prime\prime}$. 
In addition we assume that the condition {\rm (H)} holds 
for all the regular conjugate triplets. 
Then for any $k=1,\dotsc,n^{\prime\prime}$, 
$C_{R,k}$ is an $(N+2n^\prime-1)$-dimensional embedded submanifold of $\mathcal{G}{\times}X{\times}X$. 
\end{lemma}

\begin{proof}
Fix arbitrary $k=1,\dotsc,n^{\prime\prime}$, 
and pick up arbitrary $(z_0;x_0,y_0) \in C_{R,k}$. 
Suppose that $Z$ is expressed as 
$\{x^{\prime\prime}=\phi(z,x^\prime)\}$ near $(z_0,x_0)$ 
and 
$\{y^{\prime\prime}=\psi(z,y^\prime)\}$ near $(z_0,y_0)$ 
respectively. 
Let $H^\lambda(z,x^\prime,y^\prime)$ be the same as 
that of the definition of Condition (H). 
Then $H^\lambda(z_0,x_0^\prime,y_0^\prime)=0$ and 
$\operatorname{rank}\bigl(D_{z,x^\prime,y^\prime}H^{\lambda}(z_0,x_0^\prime,y_0^\prime)\bigr)=1$. 
Pick up nonzero row 
$\nabla_{z,x^\prime,y^\prime}H^{\lambda_l}(z_0,x_0^\prime,y_0^\prime)$ of 
$D_{z,x^\prime,y^\prime}H^{\lambda}(z_0,x_0^\prime,y_0^\prime)$. 
The implicit function theorem implies that 
$\{H^{\lambda_l}(z,x^\prime,y^\prime)=0\}$ is a hypersurface in 
$\mathcal{G} \times X \times X$ near $(z_0;x_0,y_0)$. 
Then $\{H^{\lambda}(z,x^\prime,y^\prime)=0\}$ is also a hypersurface in 
$\mathcal{G} \times X \times X$ near $(z_0;x_0,y_0)$ 
since $\operatorname{rank}\bigl(D_{z,x^\prime,y^\prime}H^{\lambda}(z_0,x_0^\prime,y_0^\prime)\bigr)=1$. Therefore 
%
%
$$
\{x^{\prime\prime}=\phi(z,x^\prime)\}
\cap
\{y^{\prime\prime}=\psi(z,y^\prime)\}
\cap
\{H^{\lambda}(z,x^\prime,y^\prime)=0\}
\subset 
C_{R,k}
$$
%
% 
near $(z_0;x_0,y_0)$ since 
$(z_0;x_0,y_0)$ is a regular $Z$-conjugate triplet of degree $k$. 
Using the Condition (H) again, we deduce that 
the connected component of $C_{R,k}$ containing $(z_0;x_0,y_0)$ is characterized by 
%
%
$$
F(x^{\prime\prime},y^{\prime\prime},z;x^\prime,y^\prime)
:=
\begin{bmatrix}
x^{\prime\prime}-\phi(z,x^\prime)
\\
y^{\prime\prime}-\psi(z,y^\prime) 
\\
H^\lambda(z,x^\prime,y^\prime)
\end{bmatrix}
=0
$$
%
%
near $(z_0;x_0,y_0)$. 
The differential of $F$ at $(z_0;x_0,y_0)$ is 
%
%
\begin{align*}
& DF(x_0^{\prime\prime},y_0^{\prime\prime},z_0;x_0^\prime,y_0^\prime)
\\
  =
& \begin{bmatrix}
  I_{n^{\prime\prime}} & O & -\phi_z(z_0,x_0^\prime) & -\phi_{x^\prime}(z_0,x_0^\prime) & O
  \\
  O & I_{n^{\prime\prime}} &  -\psi_z(z_0,x_0^\prime) & O & -\psi_{x^\prime}(z_0,y_0^\prime)
  \\
  O & O & H_z^\lambda(z_0,x_0^\prime,y_0^\prime) & H_{x^\prime}^\lambda(z_0,x_0^\prime,y_0^\prime) & H_{y^\prime}^\lambda(z_0,x_0^\prime,y_0^\prime)
  \end{bmatrix}. 
\end{align*}
%
%
Define an upper triangular matrix $Q$ by 
%
%
$$
Q(z,x^\prime,y^\prime)
:=
\begin{bmatrix}
I_{n^{\prime\prime}} & O & \phi_z(z,x^\prime) & \phi_{x^\prime}(z,x^\prime) & O 
\\
  & I_{n^{\prime\prime}} & \psi_z(z,y^\prime) & O & \phi_{x^\prime}(z,y^\prime) 
\\
  &                      & I_N                & O            & O
\\
  &                      &                    & I_{n^\prime} & O
\\
  &                      &                    &              & I_{n^\prime}
\end{bmatrix}.
$$
%
%
Then $\operatorname{det}\bigl(Q(z,x^\prime,y^\prime)\bigr)\equiv1$ and 
%
%
$$
DF(x_0^{\prime\prime},y_0^{\prime\prime},z_0;x_0^\prime,y_0^\prime)
Q(z_0,x_0^\prime,y_0^\prime)
=
\begin{bmatrix}
I_{2n^{\prime\prime}} & O
\\
O & D_{z,x^\prime,y^\prime}H^\lambda(z_0,x_0^\prime,y_0^\prime) 
\end{bmatrix}. 
$$
%
%
Hence $\operatorname{rank}\bigl(DF(x_0^{\prime\prime},y_0^{\prime\prime},z_0;x_0^\prime,y_0^\prime)\bigr)=2n^{\prime\prime}+1$, and we deduce that $C_{R,k}$ is an embedded 
submanifold of $\mathcal{G} \times X \times X$ of dimension 
$(N+2n)-(2n^{\prime\prime}+1)=N+2n^\prime-1$. 
\end{proof}

\end{document}
